\chapter{Higher-Order De Bruijn GNNs for Temporal Networks}

% =============================================================================
\section{Where We Stand: Three Approaches to the Same Problem}
% =============================================================================

The previous two chapters introduced two fundamentally different strategies for making Graph
Neural Networks (GNNs) respect the arrow of time in a temporal contact network.

The Saram\"aki framework (Chapter~\ref{ch:saramaki}) takes the most radical approach: it
discards the original graph entirely and rebuilds a new one, the Temporal Event Graph (TEG),
in which every contact event $(u, v, t)$ becomes a \emph{node} and causal adjacencies between
events become \emph{edges}. The result is a Directed Acyclic Graph (DAG) in which every edge is,
by construction, a valid causal step. The price is graph bloat: in the worst case the TEG has
$O(E^2)$ edges for $E$ original contact events.

The Ru et al.\ framework (Chapter~\ref{ch:ru}) keeps the original static graph but attaches a
binary \emph{temporal activation vector} of length $T$ to every edge, recording which time steps
that edge was active. No graph topology is modified. The cost is a different kind of memory: every
edge now carries a $T$-dimensional feature tensor, and the GNN must learn to read these vectors
to reason about sequences.

This chapter introduces the third and most theoretically rich approach, proposed by Qarkaxhija,
Perri, and Scholtes in \textit{``De Bruijn goes Neural: Causality-Aware Graph Neural Networks
for Time Series Data on Dynamic Graphs''} \cite{qarkaxhija2022}. Their key claim is that both
previous approaches suffer from a single root cause: they treat the \emph{first-order}
static graph---where nodes are people and edges are social ties---as the underlying message-passing
substrate. A first-order graph is inherently \emph{Markovian}: when a message arrives at node
$B$, the graph has no memory of where it came from. It will propagate to all of $B$'s neighbours
equally, regardless of which path was physically plausible.

The De Bruijn approach fixes this by operating on a completely different graph, one whose nodes
are not people but \emph{sequences of contacts}. This chapter builds up that idea from scratch.

% =============================================================================
\section{The Markovian Problem: A Concrete Failure Mode}
% =============================================================================

\begin{conceptbox}{What Does ``Markovian'' Mean?}
A process is \textbf{Markovian} (also called \emph{memoryless}) if the probability of the next
state depends \emph{only} on the current state, not on how you arrived there. A standard static
GNN is Markovian in exactly this sense. At layer $l$, node $B$ aggregates messages from all its
neighbours. It does not know---and cannot know---which of those neighbours sent the message in
layer $l-1$. Every aggregation step starts fresh with no memory of the previous step's origin.
\end{conceptbox}

Why does this matter for epidemic spreading? Consider the following example, which we will use
throughout this chapter.

\begin{examplebox}{The Causal Violation Example}
A small temporal contact network involving four people: Alice (A), Bob (B), Charlie (C):
\begin{itemize}[noitemsep]
    \item Contact $e_1$: Alice meets Bob at $t = 1$.
    \item Contact $e_2$: Bob meets Charlie at $t = 2$.
    \item Contact $e_3$: Alice meets Charlie at $t = 3$.
\end{itemize}

\textbf{Reality (causal paths):}
\begin{itemize}[noitemsep]
    \item $A \to B \to C$ is valid: events are $e_1$ then $e_2$, with $t_2 > t_1$. The virus
          can travel this route.
    \item $A \to C$ directly via $e_3$ is valid: $t_3 = 3$ and Alice was infected by $t=1$.
    \item $C \to B \to A$ is \textbf{impossible}: this would require Bob to infect Alice at
          $t=1$ after catching it from Charlie at $t=2$, which is time travel.
\end{itemize}

\textbf{What a static GNN sees:}
The static projection collapses the three events into a static triangle $A-B-C$. At message
passing layer 1, the GNN sends a message from $C$ to $B$. At layer 2, it sends the message from
$B$ to $A$. It has hallucinated the impossible path $C \to B \to A$. This is called
\textbf{causal leakage}.
\end{examplebox}

\begin{analogybox}{The News Telegraph Analogy}
Imagine information spreading in 1850 via telegraph, where messages can only travel forward in
time along existing telegraph lines.

A \emph{static map} of telegraph lines would suggest that any city can communicate with any
other city, as long as a path of lines exists. But of course, a message sent from Paris on
Monday cannot ``reply'' to a message that London sent on Tuesday---the Tuesday message does not
exist yet when the Monday message travels.

A Markovian GNN on the static map is like a cartographer who draws circles of ``reachability''
ignoring when any message was sent. The De Bruijn GNN is like a telegraph historian who carefully
records every sequence of transmissions and only connects cities that actually had a valid
chronological relay.
\end{analogybox}

\begin{mathbox}{Why DAG-GNNs and Backtracking Networks Are Still Partly Markovian}
The DAG-GNN correctly prohibits backward-in-time edges, so it eliminates the worst causal
leakage. However, within a single GNN message-passing layer, aggregation over all incoming causal
edges is still memoryless: node $B$ cannot distinguish ``I received this message because $A$
was infected first and then met me'' from ``I received this message because $C$ met me before
meeting $A$.''

The Backtracking Network's edge activation vectors encode \emph{which time steps} an edge was
active, but the node update rule still aggregates all neighbours' messages without regard for
\emph{which activation pattern led to which other activation pattern}.

The De Bruijn approach eliminates this residual Markovian assumption by making the entire
\emph{sequence} of contacts---not just individual contacts---the primary unit of representation.
\end{mathbox}

% =============================================================================
\section{What Is a De Bruijn Graph? From DNA to Sequences}
% =============================================================================

The mathematical structure underlying this paper has a long history in combinatorics and
bioinformatics. To understand it intuitively, we begin with the original application: DNA
sequence assembly.

\subsection{De Bruijn Graphs for DNA Sequences}

When biologists sequence a genome, they cannot read the entire DNA strand at once. Sequencing
machines produce millions of short fragments, each of length $k$ nucleotides. These fragments
are called \textbf{$k$-mers}. The problem of genome assembly is: given all the $k$-mers, in
what order do they originally appear?

The De Bruijn graph provides an elegant solution.

\begin{conceptbox}{De Bruijn Graphs for $k$-mers}
Let $\Sigma$ be an alphabet (for DNA: $\{A, T, G, C\}$). A \textbf{$k$-mer} is any sequence of
$k$ characters from $\Sigma$.

A \textbf{De Bruijn graph of order $k$} is constructed as follows:
\begin{itemize}[noitemsep]
    \item \textbf{Nodes:} Every distinct $(k-1)$-mer that appears in the data.
    \item \textbf{Directed edges:} A directed edge exists from node $\alpha$ to node $\beta$ if
          and only if the last $k-2$ characters of $\alpha$ match the first $k-2$ characters of
          $\beta$, AND the combined sequence $\alpha[1]\cdots\alpha[k-1]\,\beta[k-1]$ is a valid
          $k$-mer in the data.
\end{itemize}
In other words: an edge $\alpha \to \beta$ encodes the fact that the sequence $\alpha$ can be
immediately followed by character $\beta[k-1]$. Each edge in the graph represents one $k$-mer;
each node represents one $(k-1)$-mer.
\end{conceptbox}

\begin{examplebox}{De Bruijn Graph: A Small DNA Example ($k=3$)}
Suppose we observe the following 3-mers from a DNA fragment: \texttt{ACG}, \texttt{CGT},
\texttt{GTG}, \texttt{TGA}.

The nodes (2-mers) are: \texttt{AC}, \texttt{CG}, \texttt{GT}, \texttt{TG}, \texttt{GA}.

The edges (representing 3-mers):
\begin{itemize}[noitemsep]
    \item \texttt{ACG}: node \texttt{AC} $\to$ node \texttt{CG} \hfill (overlap: \texttt{C})
    \item \texttt{CGT}: node \texttt{CG} $\to$ node \texttt{GT} \hfill (overlap: \texttt{G})
    \item \texttt{GTG}: node \texttt{GT} $\to$ node \texttt{TG} \hfill (overlap: \texttt{T})
    \item \texttt{TGA}: node \texttt{TG} $\to$ node \texttt{GA} \hfill (overlap: \texttt{G})
\end{itemize}
Following the path \texttt{AC} $\to$ \texttt{CG} $\to$ \texttt{GT} $\to$ \texttt{TG} $\to$
\texttt{GA} in the De Bruijn graph spells out the original sequence \texttt{ACGTGA}.
\end{examplebox}

The key insight is the \textbf{overlap rule}: an edge $\alpha \to \beta$ exists only when the
\emph{suffix} of $\alpha$ matches the \emph{prefix} of $\beta$. This overlap constraint ensures
that consecutive edges form a valid, non-contradictory sequence. Genome assembly reduces to
finding an Eulerian path through this graph.

\subsection{Adapting the Concept to Temporal Networks}

Now we replace the DNA alphabet with a contact network. Instead of characters $\{A, T, G, C\}$,
our ``alphabet'' is the set of all possible contacts $(u, v, t)$ between people. Instead of
$k$-mers (sequences of $k$ characters), we have \textbf{temporal walks of length $k$}: sequences
of $k$ contacts that are temporally ordered and share intermediate nodes.

\begin{conceptbox}{Temporal Walk of Length $k$}
A \textbf{temporal walk of length $k$} in a dynamic graph $G^T = (V, E^T)$ is a sequence of
$k$ contact events:
\[ w = (e_{i_1},\, e_{i_2},\, \ldots,\, e_{i_k}) \]
where each $e_{i_j} = (u_j, v_j, t_j) \in E^T$, subject to two conditions:
\begin{enumerate}[noitemsep]
    \item \textbf{Node continuity:} The destination of one event equals the source of the next.
          Formally, $v_j = u_{j+1}$ for all $j \in \{1, \ldots, k-1\}$.
    \item \textbf{Temporal ordering:} Events occur in strictly increasing time.
          Formally, $t_1 < t_2 < \cdots < t_k$ (with optional maximum gap $\delta$:
          $t_{j+1} - t_j \leq \delta$).
\end{enumerate}
Such a walk represents a \emph{physically possible} transmission path of length $k$.
\end{conceptbox}

The concept directly parallels the DNA case: the ``alphabet'' is the set of directed contact
events; the ``$k$-mers'' are valid temporal walks of length $k$; and the ``overlap rule'' is
the combined requirement of shared intermediate node and strictly increasing timestamps.

% =============================================================================
\section{Constructing De Bruijn Graphs for Temporal Networks}
% =============================================================================

We now define the central mathematical object of this chapter.

\subsection{The $k$-th Order De Bruijn Graph}

\begin{mathbox}{Definition: $k$-th Order De Bruijn Graph $G^{(k)}$}
Let $G^T = (V, E^T)$ be a dynamic graph with node set $V$ and time-stamped directed edges
$E^T \subseteq V \times V \times \mathbb{N}$.

The \textbf{$k$-th order De Bruijn graph} of $G^T$ is a static directed graph
$G^{(k)} = (V^{(k)}, E^{(k)})$ where:

\medskip
\textbf{Node set $V^{(k)}$:} Every valid temporal walk of length $k-1$ in $G^T$ is one node.
Formally, $u \in V^{(k)}$ is a tuple $u = (u_0, u_1, \ldots, u_{k-1})$ of $k$ nodes from $V$
such that there exists a causal walk visiting $u_0, u_1, \ldots, u_{k-1}$ in order.

\medskip
\textbf{Edge set $E^{(k)}$:} A directed edge $(u, v) \in E^{(k)}$ exists if and only if:
\begin{enumerate}[noitemsep]
    \item $u = (u_0, \ldots, u_{k-1})$ and $v = (v_0, \ldots, v_{k-1})$ overlap: $u_i = v_{i-1}$
          for $i = 1, \ldots, k-1$.
    \item The combined walk $(u_0, u_1, \ldots, u_{k-1}, v_{k-1})$ is a valid causal walk of
          length $k$ in $G^T$.
\end{enumerate}
Each edge $(u, v)$ in $G^{(k)}$ thus represents a valid temporal walk of length $k$.
\end{mathbox}

\begin{keyinsight}{What the Order $k$ Controls}
The order $k$ determines the \emph{memory depth} of the model:
\begin{itemize}[noitemsep]
    \item $k = 1$: Nodes represent individual contacts. Edges represent pairs of causally
          adjacent contacts. This is approximately the same as the Temporal Event Graph (TEG)
          from the Saram\"aki chapter, without the DAG structure.
    \item $k = 2$: Nodes represent \emph{pairs} of consecutive contacts
          (two-step transmission sequences). Edges connect sequences that share a one-step
          overlap. The GNN now passes messages between two-step transmission sequences.
    \item $k = 3$: Nodes represent three-step sequences. The model ``remembers'' three hops.
\end{itemize}
A first-order static graph is a zeroth-order approximation: it records that contacts exist, but
has no memory of sequences at all.
\end{keyinsight}

\subsection{Our Implementation: First-Order De Bruijn Graph ($k=1$)}

In the implementation used in this thesis (see \texttt{utils/make\_de\_bruijn\_graph.py} and
\texttt{gnn/graph\_builder.py}), we use a first-order De Bruijn graph, which corresponds to the
case $k=1$ in the paper's notation. This means:
\begin{itemize}[noitemsep]
    \item \textbf{Nodes:} Every directed contact event $(u, v, t)$ in the original network
          becomes a node in the De Bruijn graph.
    \item \textbf{Directed edges:} An edge $(u, v, t_1) \to (x, y, t_2)$ exists iff $v = x$
          (shared bridge node) and $t_2 > t_1$ (time ordering).
\end{itemize}
This is the representation used in all code walkthroughs below. The higher-order variants ($k>1$)
follow the same principle but with nodes representing multi-hop sequences.

% =============================================================================
\section{Worked Example: Building the De Bruijn Graph Step by Step}
% =============================================================================

We now trace through the complete construction of a first-order De Bruijn graph on a concrete
small example, so that every step is visible.

\subsection{The Input Temporal Network}

Consider a temporal network with four people: Alice (A), Bob (B), Charlie (C), Dana (D). There
are four contact events:
\begin{align*}
    e_1 &= (A, B, t=1) \\
    e_2 &= (B, C, t=2) \\
    e_3 &= (A, C, t=3) \\
    e_4 &= (C, D, t=4)
\end{align*}

In the \emph{static projection} (ignoring timestamps), this looks like a graph with edges
$A-B$, $B-C$, $A-C$, $C-D$.

\subsection{Step 1: Enumerate De Bruijn Nodes}

Each directed version of each contact event becomes one De Bruijn node. Because the original
contacts are undirected (a meeting between $A$ and $B$ is the same whether we write it as
$(A \to B)$ or $(B \to A)$), each contact event generates \emph{two} De Bruijn nodes:

\begin{center}
\begin{tabular}{ll}
\toprule
\textbf{Original event} & \textbf{De Bruijn nodes created} \\
\midrule
$e_1 = (A, B, t=1)$ & $(A, B, 1)$ and $(B, A, 1)$ \\
$e_2 = (B, C, t=2)$ & $(B, C, 2)$ and $(C, B, 2)$ \\
$e_3 = (A, C, t=3)$ & $(A, C, 3)$ and $(C, A, 3)$ \\
$e_4 = (C, D, t=4)$ & $(C, D, 4)$ and $(D, C, 4)$ \\
\bottomrule
\end{tabular}
\end{center}

In addition, the builder adds \textbf{sentinel nodes} $(n, n, t_{\min})$ and
$(n, n, t_{\max})$ for each person $n$. We will explain sentinels in detail in
Section~\ref{sec:sentinels}. For now, note that each node is a triple $(u, v, t)$ where $u$
is the ``source'' person, $v$ is the ``destination'' person, and $t$ is the contact time.

\subsection{Step 2: Draw De Bruijn Edges}

For each pair of De Bruijn nodes $(u, v, t_1)$ and $(x, y, t_2)$, we draw a directed edge
$(u, v, t_1) \to (x, y, t_2)$ if and only if $v = x$ (the destination of the first event
equals the source of the second) \emph{and} $t_2 > t_1$.

Let us check all valid pairs:

\begin{examplebox}{Checking Every Candidate Edge}
\begin{itemize}[noitemsep]
    \item $(A,B,1) \to (B,C,2)$: bridge node $B=B$ \checkmark, time $2>1$ \checkmark.
          \textbf{Edge exists.} This encodes the causal path $A \to B \to C$.
    \item $(A,B,1) \to (B,A,1)$: bridge $B=B$ \checkmark, time $1>1$? \texttimes. No edge.
          (Same timestamp, not strictly later.)
    \item $(B,C,2) \to (C,D,4)$: bridge $C=C$ \checkmark, time $4>2$ \checkmark.
          \textbf{Edge exists.} Encodes $B \to C \to D$.
    \item $(B,C,2) \to (C,A,3)$: bridge $C=C$ \checkmark, time $3>2$ \checkmark.
          \textbf{Edge exists.} Encodes $B \to C \to A$ (backwards reach, but causally valid).
    \item $(A,C,3) \to (C,D,4)$: bridge $C=C$ \checkmark, time $4>3$ \checkmark.
          \textbf{Edge exists.} Encodes $A \to C \to D$.
    \item $(A,B,1) \to (A,C,3)$: bridge $B \neq A$. No edge. Alice did not pass through Bob
          to reach Charlie directly.
    \item $(C,B,2) \to (B,A,1)$: bridge $B=B$ \checkmark, but time $1 > 2$? \texttimes.
          No edge. Time moves backwards here; causally impossible.
\end{itemize}
The complete set of forward causal edges is thus:
$(A,B,1) \to (B,C,2)$, $(B,C,2) \to (C,D,4)$, $(B,C,2) \to (C,A,3)$, $(A,C,3) \to (C,D,4)$.
\end{examplebox}

\subsection{Step 3: Visualise the Result}

\begin{figure}[H]
\centering
\begin{tikzpicture}[
    evtnode/.style={
        rectangle, rounded corners=4pt, draw=orange!80!black,
        fill=orange!12, thick, minimum width=1.6cm, minimum height=0.7cm,
        font=\footnotesize\bfseries, align=center
    },
    sentnode/.style={
        circle, draw=blue!70!black, fill=blue!12, thick,
        minimum size=0.8cm, font=\scriptsize\bfseries, align=center
    },
    fwd/.style={-{Stealth[length=2.5mm]}, thick, orange!80!black},
    sent/.style={-{Stealth[length=2mm]}, thick, blue!60, dashed}
]
% Event nodes (top row, left to right in time)
\node[evtnode] (AB) at (0,  0)   {$A{\to}B$\\$t=1$};
\node[evtnode] (BA) at (0, -2.2) {$B{\to}A$\\$t=1$};

\node[evtnode] (BC) at (3.5,  0)   {$B{\to}C$\\$t=2$};
\node[evtnode] (CB) at (3.5, -2.2) {$C{\to}B$\\$t=2$};

\node[evtnode] (AC) at (7,  0)   {$A{\to}C$\\$t=3$};
\node[evtnode] (CA) at (7, -2.2) {$C{\to}A$\\$t=3$};

\node[evtnode] (CD) at (10.5,  0)   {$C{\to}D$\\$t=4$};
\node[evtnode] (DC) at (10.5, -2.2) {$D{\to}C$\\$t=4$};

% Causal De Bruijn edges
\draw[fwd] (AB) -- node[above, font=\scriptsize]{B=B} (BC);
\draw[fwd] (BC) -- node[above, font=\scriptsize]{C=C} (CD);
\draw[fwd] (BC) to[bend right=18] node[below, font=\scriptsize]{C=C} (CA);
\draw[fwd] (AC) -- node[above, font=\scriptsize]{C=C} (CD);

% Sentinel nodes along the bottom
\node[sentnode] (sA) at (-1.8, -4.5) {\small$A{:}t_0$};
\node[sentnode] (sB) at (1.8, -4.5)  {\small$B{:}t_0$};
\node[sentnode] (sC) at (5.3, -4.5)  {\small$C{:}t_0$};
\node[sentnode] (sD) at (8.8, -4.5)  {\small$D{:}t_0$};

% Sentinel start-edges (from person's sentinel to their first event)
\draw[sent] (sA) to[bend left=10] (AB);
\draw[sent] (sA) to[bend right=10] (AC);
\draw[sent] (sB) to[bend left=10] (BA);
\draw[sent] (sB) to[bend right=10] (BC);
\draw[sent] (sC) to[bend left=10] (CB);
\draw[sent] (sC) to[bend right=10] (CA);
\draw[sent] (sC) -- (CD);
\draw[sent] (sD) -- (DC);
\end{tikzpicture}
\caption{First-order De Bruijn graph for the example temporal network with four contact events.
Orange rectangles are \emph{event nodes}; blue circles are \emph{sentinel nodes} (one per person).
Orange arrows are causal De Bruijn edges; blue dashed arrows connect each sentinel to the events
it ``owns.'' Only forward-in-time causal edges are shown for clarity.}
\label{fig:debruijn_example}
\end{figure}

Notice what the graph has achieved. The directed path
$(A,B,1) \to (B,C,2) \to (C,D,4)$
in the De Bruijn graph precisely encodes the causal chain
$A$ infects $B$ at time 1, then $B$ infects $C$ at time 2, then $C$ infects $D$ at time 4.
There is \emph{no edge} connecting $(C,B,2)$ to $(B,A,1)$, because that would require time
to flow backward. The impossible path $C \to B \to A$ cannot be traversed.

\begin{keyinsight}{Causal Leakage is Structurally Impossible}
In the De Bruijn graph, every directed path is, by construction, a valid causal transmission
chain. There is no need to check time orderings during message passing because all invalid paths
have already been removed during graph construction. Causal correctness is encoded in the
\emph{topology}, not the \emph{features}.

This is the central advantage over both the DAG-GNN approach (which also removes impossible edges
but at higher construction cost for large $k$) and the Backtracking Network (which encodes time
in features, leaving the GNN to learn which sequences are valid).
\end{keyinsight}

% =============================================================================
\section{Why Higher-Order Graphs Capture Non-Markovian Dynamics}
% =============================================================================

\subsection{First-Order versus Second-Order: An Epidemic Illustration}

Consider a temporal network representing a hospital ward. Staff members can spread a pathogen
by touching contaminated surfaces. Suppose we observe the following pattern in the contact data:
\begin{itemize}[noitemsep]
    \item Doctors often visit Room 101 immediately \emph{after} attending the morning handover
          meeting in Room 100.
    \item After visiting Room 101, they typically go to Room 102.
\end{itemize}

A \emph{first-order} De Bruijn graph records: ``Doctor D contacted Room 100. Doctor D contacted
Room 101. Doctor D contacted Room 102.'' Each contact is a node; edges connect contacts that
share a person.

A \emph{second-order} De Bruijn graph records, as single nodes, the \emph{two-step sequences}:
``D went from Room 100 to Room 101'' and ``D went from Room 101 to Room 102.'' The edge between
these two-step nodes encodes the full three-step sequence ``Room 100 $\to$ Room 101 $\to$ Room
102.'' This is the \emph{pathway}, the transmission chain that the pathogen actually travels.

\begin{analogybox}{The Train Timetable Analogy}
Think of a rail network. A first-order graph knows which cities are connected by rail. A
second-order graph knows which \emph{connections} are available: ``Passengers who arrive at
Munich from Frankfurt can then board the Munich-to-Vienna train.'' This is crucial information
if you want to know whether a traveller can get from Frankfurt to Vienna in one day: not just
whether the rail links exist, but whether the connection times work out.

Epidemic spreading obeys the same logic. A pathogen is like a traveller. It cannot ``take a
connection'' if the first train arrives after the second one has departed. Higher-order De Bruijn
graphs are timetable-aware; first-order graphs are not.
\end{analogybox}

\subsection{Mathematical Statement: Non-Markovian Message Passing}

In a standard GNN operating on $G^{(1)}$ (the first-order static graph), the update rule at
layer $l$ is:
\[
    \mathbf{h}_v^{(l)} = \sigma\!\left(\mathbf{W}^{(l)} \cdot \textsc{Agg}\!\left(
        \left\{\mathbf{h}_u^{(l-1)} : (u, v) \in E \right\}\right)\right)
\]
Node $v$ has no information about \emph{which} of its neighbours $u$ sent the message in the
previous layer, so it cannot distinguish between a causal path $A \to v \to w$ and a
non-causal one.

When the GNN operates on $G^{(k)}$ (the $k$-th order De Bruijn graph), each node
$\tilde{u} = (u_0, \ldots, u_{k-1}) \in V^{(k)}$ represents an entire $k$-step walk. The
update rule becomes:

\begin{equation}
    \bar{h}_{v}^{k,l} = \sigma\!\left(\mathbf{W}^{k,l} \cdot
    \sum_{\substack{u \in V^{(k)}: \\ (u,v) \in E^{(k)}}}
    \frac{w(u,v) \cdot \bar{h}_u^{k,l-1}}{\sqrt{S(v) \cdot S(u)}}\right)
    \label{eq:dbgnn_update}
\end{equation}

where $w(u,v)$ is the edge weight (encoding how frequently the two-step sequence was observed),
$\mathbf{W}^{k,l}$ is a trainable weight matrix, $S(v) = \sum_{u} w(u,v)$ is the weighted
degree, and $\sigma$ is a nonlinear activation. Because $v$ represents a $k$-step walk, the
message $\bar{h}_u^{k,l-1}$ already encodes the full $(k-1)$-step history of the walk leading
to $v$. The aggregation is therefore \emph{non-Markovian}: the ``history'' is built into the
node identity itself.

% =============================================================================
\section{The Sentinel Node: Mapping Back to People}
\label{sec:sentinels}
% =============================================================================

After building the De Bruijn graph and running GNN message passing on it, we face a structural
mismatch:

\begin{itemize}
    \item Our \textbf{labels} are defined over \emph{people} (original nodes $v \in V$).
    \item The \textbf{De Bruijn graph nodes} are contact events $(u, v, t)$, not people.
\end{itemize}

A person $A$ who participated in three contacts now ``lives'' across three distinct De Bruijn
nodes: $(A, B, 1)$, $(B, A, 1)$, $(A, C, 3)$. How do we extract a single
prediction score for person $A$?

\begin{analogybox}{The Personal Logbook Analogy}
Imagine every person in the network carries a personal logbook. Every time they have a contact,
they write an entry in the logbook. At the end of the observation period, the logbook contains a
chronological summary of all that person's interactions.

The \textbf{sentinel node} $(v, v, t_{\text{end}})$ is exactly this logbook. It is a special
De Bruijn node that represents ``person $v$'s record at the end of the observation window.''
Every contact event involving $v$ sends information \emph{into} the sentinel (via the De Bruijn
edges from event nodes to the sentinel). After message passing, the sentinel's embedding is a
summary of all the causal information that flowed through person $v$ over time.

When we want to score person $v$ as a potential epidemic source, we read the final entry in
their logbook: we extract the embedding of their sentinel node $(v, v, t_{\text{end}})$.
\end{analogybox}

\subsection{Formal Definition of Sentinel Nodes}

\begin{mathbox}{Sentinel Nodes}
For each person $v \in V$ in the original network, we add two \textbf{sentinel nodes} to the De
Bruijn graph:
\begin{itemize}[noitemsep]
    \item $(v, v, t_{\text{start}})$: the \textbf{start sentinel}. Connected forward to every
          event node $(v, w, t)$ for $t > t_{\text{start}}$, encoding ``at the start of the
          observation, person $v$ had not yet been involved in any contact.''
    \item $(v, v, t_{\text{end}})$: the \textbf{end sentinel}. Every event node $(u, v, t)$
          or $(v, u, t)$ connects forward to this sentinel. It is the \emph{readout anchor}:
          after message passing, we extract the embedding of this node to obtain the
          source-detection score for person $v$.
\end{itemize}
An edge $(v, v, t_{\text{start}}) \to (v, v, t_{\text{end}})$ is also added to ensure the
sentinel is reachable even for people with no contacts.
\end{mathbox}

\begin{examplebox}{Sentinel Edges in Our Running Example}
For person Alice (A), with $t_{\text{start}} = 1$ and $t_{\text{end}} = 4$:
\begin{itemize}[noitemsep]
    \item Start sentinel $(A,A,1)$ connects forward to event nodes $(A,B,1)$ and $(A,C,3)$
          (contacts where Alice is the source).
    \item Event nodes $(A,B,1)$ and $(A,C,3)$ connect to end sentinel $(A,A,4)$.
    \item After GNN message passing, the embedding $\mathbf{h}_{(A,A,4)}$ contains aggregated
          information from all paths flowing through Alice.
    \item This embedding is used to score Alice as a potential source.
\end{itemize}
\end{examplebox}

% =============================================================================
\section{The DBGNN Architecture: Putting It All Together}
% =============================================================================

With the De Bruijn graph constructed and sentinel nodes defined, the DBGNN architecture
proceeds in four clean stages.

\begin{conceptbox}{DBGNN Pipeline Overview}
\begin{enumerate}
    \item \textbf{Feature initialisation:} Map original SIR node features to De Bruijn node
          embeddings. Event nodes and sentinel nodes receive different projections because they
          encode different information (one contact vs.\ one person's summary).
    \item \textbf{De Bruijn message passing:} Run $L$ layers of GNN message passing on the De
          Bruijn graph. Messages flow only along valid causal edges.
    \item \textbf{Sentinel readout:} Extract the embedding of the end-time sentinel
          $(v, v, t_{\text{end}})$ for each person $v$. This is the person-level representation.
    \item \textbf{Score projection and masking:} Apply a linear layer to each sentinel embedding
          to produce a scalar score. Mask out nodes known to be susceptible (they cannot be
          the source). Apply log-softmax to obtain a probability distribution.
\end{enumerate}
\end{conceptbox}

\subsection{Stage 1: Feature Initialisation}

Every De Bruijn node $(u, v, t)$ corresponds to a pair of original nodes $(u, v)$. We can look
up the SIR state vectors of both:
\[
    \mathbf{x}_u \in \{0,1\}^3, \quad \mathbf{x}_v \in \{0,1\}^3
\]
For \textbf{event nodes} ($u \neq v$), the initial embedding is a projection of the
\emph{concatenation} of both endpoints' features, giving a 6-dimensional input:
\[
    \mathbf{h}_{\text{event}} = \text{ReLU}\!\left(\mathbf{W}_{\text{event}}\, [\mathbf{x}_u \| \mathbf{x}_v]\right)
    \in \mathbb{R}^D
\]
For \textbf{sentinel nodes} ($u = v$), only one person is involved, so the input is 3-dimensional:
\[
    \mathbf{h}_{\text{sentinel}} = \text{ReLU}\!\left(\mathbf{W}_{\text{sentinel}}\, \mathbf{x}_u\right)
    \in \mathbb{R}^D
\]
The correct projection is selected for each node based on the boolean mask \texttt{is\_sentinel}.

\subsection{Stage 2: De Bruijn Message Passing}

The layer update is a mean-aggregator SAGE-style step:
\[
    \mathbf{h}_{v}^{(l+1)} = \text{ReLU}\!\left(
        \mathbf{W}_{\text{self}}\, \mathbf{h}_{v}^{(l)}
        + \mathbf{W}_{\text{agg}}\, \frac{1}{\deg^{-}(v)} \sum_{u: (u,v) \in E_{\text{DB}}}
          \mathbf{h}_{u}^{(l)}
    \right)
\]
where $\deg^{-}(v)$ is the in-degree of node $v$ in the De Bruijn graph (the number of causal
predecessors). The two linear transforms $\mathbf{W}_{\text{self}}$ and $\mathbf{W}_{\text{agg}}$
are learnable parameters, each of size $D \times D$.

\subsection{Stage 3: Sentinel Readout}

After $L$ message-passing layers, information has propagated forward along all valid causal chains.
For each original person $v$, we extract the embedding of the end-time sentinel:
\[
    \mathbf{e}_v = \mathbf{h}_{(v,v,t_{\text{end}})}^{(L)} \in \mathbb{R}^D
\]
This gives a matrix $\mathbf{E} \in \mathbb{R}^{N \times D}$ of person-level representations,
where $N = |V|$ is the number of people.

\subsection{Stage 4: Score Projection and Log-Softmax}

A final linear layer maps each person's embedding to a scalar logit:
\[
    s_v = \mathbf{w}_{\text{out}}^T \mathbf{e}_v \in \mathbb{R}
\]
Nodes whose SIR state is Susceptible are set to $-\infty$ (they cannot be the source since they
were never infected). A log-softmax is then applied over all non-masked nodes:
\[
    \hat{p}_v = \log\frac{\exp(s_v)}{\sum_{w: \text{non-susceptible}} \exp(s_w)}
\]
The output $\hat{\mathbf{p}} \in \mathbb{R}^N$ is a vector of log-probabilities, one per person.
Training uses the negative log-likelihood of the true source.

% =============================================================================
\section{Code Walkthrough: \texttt{gnn/dbgnn.py} and \texttt{gnn/graph\_builder.py}}
% =============================================================================

We now trace through the actual implementation, connecting each architectural stage to the
corresponding code in the codebase.

\subsection{Graph Construction: \texttt{utils/make\_de\_bruijn\_graph.py}}

The De Bruijn graph is constructed by the function \texttt{make\_de\_bruijn\_graph} in
\texttt{utils/make\_de\_bruijn\_graph.py}. Its logic follows the three steps of the previous
sections exactly.

\begin{codewalk}{Building the De Bruijn Graph Node-by-Node}
The function receives \texttt{H\_array}, a NumPy array of shape $[E, 3]$ where each row is a
contact event $(u, v, t)$, sorted by time.

\begin{lstlisting}[language=Python]
def make_de_bruijn_graph(H_array, start_t, end_t,
                          time_reverse=False, directed=False):
    B = nx.DiGraph()
    nodes = set(H_array[:, 0]).union(H_array[:, 1])

    # --- Sentinel nodes for every person ---
    for n in nodes:
        B.add_node((n, n, end_t))    # end-time sentinel
        B.add_node((n, n, start_t))  # start-time sentinel
    for n in nodes:
        # Direct edge from start to end sentinel (base case)
        B.add_edge((n, n, start_t), (n, n, end_t), ...)
\end{lstlisting}

The loop then processes contacts chronologically. For each contact $(u, v, t)$:
\begin{lstlisting}[language=Python]
    for u, v, t in H_array:
        # Connect previous events involving u to this new event
        for x, y, t1 in previous_contacts[u]:
            B.add_edge((x, y, t1), (u, v, t), ...)
        # Connect start sentinel to this event
        B.add_edge((u, u, start_t), (u, v, t), ...)
        # Connect this event to end sentinel
        B.add_edge((u, v, t), (v, v, end_t), ...)
        # Record event for future use
        previous_contacts[v].append((u, v, t))
\end{lstlisting}
The variable \texttt{previous\_contacts[v]} tracks all prior contacts whose endpoint was $v$,
enabling efficient look-up of all valid predecessors when processing new contacts involving $v$.
\end{codewalk}

\begin{keyinsight}{Efficient Causal Edge Construction}
The algorithm processes contacts in sorted temporal order. For each new contact $(u, v, t)$, it
only needs to consult \texttt{previous\_contacts[u]}: the list of prior events that ended at
node $u$ (making $u$ the bridge node). This avoids checking all $O(E^2)$ pairs, reducing
construction to $O(E \cdot \bar{d})$ where $\bar{d}$ is the average number of prior contacts
per node.
\end{keyinsight}

\subsection{Tensor Preparation: \texttt{gnn/graph\_builder.py}}

The function \texttt{build\_de\_bruijn\_graph} in \texttt{gnn/graph\_builder.py} converts the
NetworkX De Bruijn graph produced above into PyTorch tensors that can be passed directly to the
DBGNN model.

\begin{codewalk}{Converting the De Bruijn Graph to Tensors}
\begin{lstlisting}[language=Python]
def build_de_bruijn_graph(H: nx.Graph, directed: bool = False) -> dict:
    # Build NetworkX De Bruijn graph from H_array
    B = _make_db(H_array, start_t, end_t, ...)

    # Assign integer index to every De Bruijn node
    node_list   = list(B.nodes())
    node_to_idx = {node: i for i, node in enumerate(node_list)}

    # Edge index tensor [2, E_db]
    edges = [(node_to_idx[u], node_to_idx[v]) for u, v in B.edges()]
    edge_index = torch.tensor(edges, dtype=torch.long).T

    # For each DB node (x, y, t): store (x, y) as original node pair
    db_node_to_original = torch.zeros(db_n_nodes, 2, dtype=torch.long)
    is_sentinel = torch.zeros(db_n_nodes, dtype=torch.bool)
    for i, node in enumerate(node_list):
        x, y, _ = node
        db_node_to_original[i, 0] = int(x)  # source person
        db_node_to_original[i, 1] = int(y)  # dest person
        if int(x) == int(y):
            is_sentinel[i] = True            # x==y means sentinel

    # For readout: find DB index of end-sentinel for each orig node
    sentinel_end_indices = torch.zeros(n_nodes, dtype=torch.long)
    for n in range(n_nodes):
        end_sentinel = (n, n, end_t)
        if end_sentinel in node_to_idx:
            sentinel_end_indices[n] = node_to_idx[end_sentinel]
\end{lstlisting}

This function returns a dictionary with five tensors:
\begin{itemize}[noitemsep]
    \item \texttt{edge\_index}: the De Bruijn graph's edge connectivity, shape $[2, E_{\text{DB}}]$.
    \item \texttt{db\_node\_to\_original}: maps each DB node to its original $(u, v)$ pair,
          shape $[N_{\text{DB}}, 2]$.
    \item \texttt{sentinel\_end\_indices}: the DB node index of the end-sentinel for each
          original person, shape $[N]$.
    \item \texttt{is\_sentinel}: boolean mask indicating which DB nodes are sentinels,
          shape $[N_{\text{DB}}]$.
\end{itemize}
\end{codewalk}

\subsection{The DBGNN Model: \texttt{gnn/dbgnn.py}}

The full model is implemented in \texttt{gnn/dbgnn.py}. We walk through its three main
components.

\begin{codewalk}{Stage 1: Feature Initialisation (Split by Node Type)}
The constructor defines two separate linear projections, one for event nodes (6-dim input)
and one for sentinel nodes (3-dim input):
\begin{lstlisting}[language=Python]
class DBGNN(torch.nn.Module):
    def __init__(self, hidden_channels, num_conv_layers, ...):
        self.proj_sentinel = Linear(3, hidden_channels)  # 3 = |SIR|
        self.proj_event    = Linear(6, hidden_channels)  # 6 = |SIR|*2
        self.convs = ModuleList(
            [DBGNNLayer(hidden_channels) for _ in range(num_conv_layers)]
        )
        self.out = Linear(hidden_channels, 1)
\end{lstlisting}

In the forward pass, both projections are computed for \emph{all} DB nodes, then we select
the right one using \texttt{torch.where} with the \texttt{is\_sentinel} mask:
\begin{lstlisting}[language=Python]
    # x: [B, N, 3]  original SIR states for all people
    x_u = x[:, db_node_to_original[:, 0], :]   # [B, db_N, 3]
    x_v = x[:, db_node_to_original[:, 1], :]   # [B, db_N, 3]

    h_sent  = F.relu(self.proj_sentinel(x_u))
    h_event = F.relu(self.proj_event(torch.cat([x_u, x_v], dim=-1)))

    # Blend: pick projection based on node type
    mask = is_sentinel.view(1, db_N, 1).expand(B, db_N, D)
    h = torch.where(mask, h_sent, h_event)     # [B, db_N, D]
\end{lstlisting}
The indexing trick \texttt{x[:, db\_node\_to\_original[:, 0], :]} simultaneously gathers the
SIR feature vectors of the source-person for \emph{all} DB nodes at once, without a Python loop.
\end{codewalk}

\begin{codewalk}{Stage 2: DBGNNLayer --- Manual Mean-Aggregation SAGE}
The \texttt{DBGNNLayer} class implements the message-passing update manually, without PyTorch
Geometric's sparse batching, to support the \texttt{[B, N, D]} batched tensor layout used
throughout this codebase.
\begin{lstlisting}[language=Python]
class DBGNNLayer(torch.nn.Module):
    def __init__(self, hidden_dim):
        self.W_self = Linear(hidden_dim, hidden_dim, bias=False)
        self.W_agg  = Linear(hidden_dim, hidden_dim, bias=True)

    def forward(self, h, edge_index):
        # h: [B, db_N, D]
        src, dst = edge_index[0], edge_index[1]

        # Step 1: gather source features
        h_src = h[:, src, :]               # [B, E_db, D]

        # Step 2: scatter-add into aggregation buffer
        dst_idx = dst.view(1, -1, 1).expand(B, E, D)
        agg_sum = h.new_zeros(B, db_N, D)
        agg_sum.scatter_add_(1, dst_idx, h_src)

        # Step 3: normalise by in-degree
        deg = h.new_zeros(db_N)
        deg.scatter_add_(0, dst, torch.ones(E, device=h.device))
        deg = deg.clamp(min=1).view(1, db_N, 1)
        agg_mean = agg_sum / deg

        # Step 4: SAGE update (self + aggregated neighbour)
        h_new = F.relu(self.W_self(h) + self.W_agg(agg_mean))
        return h_new
\end{lstlisting}
The \texttt{scatter\_add\_} operation is the vectorised equivalent of the loop
``for each edge $(u,v)$: \texttt{agg[v] += h[u]},'' computed simultaneously over all edges
and all batch elements. The \texttt{clamp(min=1)} prevents division by zero for isolated nodes.
\end{codewalk}

\begin{codewalk}{Stages 3 and 4: Sentinel Readout, Masking, and Log-Softmax}
After $L$ message-passing layers, we extract the end-sentinel embeddings and produce the final
output:
\begin{lstlisting}[language=Python]
    # Stage 3: Readout via end-time sentinels
    sent_idx = sentinel_end_indices.to(x.device)  # [N]
    node_emb = h[:, sent_idx, :]                   # [B, N, D]

    # Stage 4: Score -> mask susceptibles -> log-softmax
    scores = self.out(node_emb).squeeze(-1)         # [B, N]
    susceptible_mask = x[..., 0].bool()             # [B, N] True if S
    scores = scores.masked_fill(susceptible_mask, float("-inf"))
    return F.log_softmax(scores, dim=-1)            # [B, N]
\end{lstlisting}
The indexing \texttt{h[:, sent\_idx, :]} selects, for every person $v$ in the original network,
the embedding of that person's end-sentinel from the full De Bruijn embedding matrix. The
susceptible mask ensures that nodes in state S are assigned probability $-\infty$ before
normalisation, because a node that was never infected cannot be the source.
\end{codewalk}

% =============================================================================
\section{Model Selection: Choosing the Right Order $k$}
% =============================================================================

A practical challenge with De Bruijn graphs is that different orders $k$ represent the same
data in different ways. A higher $k$ captures more complex temporal dependencies but also
increases the graph size and the risk of overfitting to spurious patterns. How do we choose $k$?

Qarkaxhija et al.\ address this with \textbf{statistical model selection}: they compute, for
each candidate order $k$, the statistics of causal walks of length $k$ in the data and apply
a likelihood-ratio test to determine whether a $k$-th order model explains significantly more
of the data than a $(k-1)$-th order model. The optimal $k$ is the lowest order for which the
test cannot be rejected.

\begin{conceptbox}{Model Order Selection (Non-Technical Summary)}
\begin{enumerate}[noitemsep]
    \item Count how often each $k$-step sequence appears in the data (the empirical distribution
          of causal walks of length $k$).
    \item Compare this to what we would expect under a lower-order (more Markovian) model.
    \item If the observed frequencies deviate significantly from the lower-order prediction,
          increase $k$. If not, stop.
\end{enumerate}
The test is parameter-free: it requires no cross-validation, no learning. The optimal graph
order is inferred directly from the temporal statistics of the data.
\end{conceptbox}

In the empirical datasets evaluated in the paper (high-school proximity networks, workplace
contacts, hospital ward data), a second-order De Bruijn graph ($k=2$) consistently provided
significantly better explanatory power than the first-order model, confirming that non-Markovian
patterns are present and meaningful in real contact data.

% =============================================================================
\section{Comparison of the Three Temporal GNN Architectures}
% =============================================================================

\begin{table}[H]
\centering
\renewcommand{\arraystretch}{1.4}
\begin{tabular}{p{2.6cm} p{3.2cm} p{2.8cm} p{2.8cm} p{3.2cm}}
\toprule
\textbf{Property} &
\textbf{Saram{\"a}ki DAG-GNN} &
\textbf{Ru et al.\ Backtracking} &
\textbf{Qarkaxhija De Bruijn} \\
\midrule
\textbf{How time is encoded} &
Events become nodes in a DAG; edges are causal adjacencies. &
Original graph kept; edges carry binary activation vectors of length $T$. &
Contact sequences become nodes; edges connect overlapping sequences. \\
\midrule
\textbf{Graph size} &
$N_{\text{nodes}} = E$ (events). $N_{\text{edges}} \leq O(E^2)$. Potentially very large. &
Identical to static graph: $N_{\text{nodes}} = |V|$, $N_{\text{edges}} = |E|$. Memory $O(|E| \cdot T)$. &
$N_{\text{nodes}} = 2E + 2|V|$ (events + sentinels). $N_{\text{edges}} \leq O(E \cdot \bar{d})$. \\
\midrule
\textbf{Temporal resolution} &
Exact: every event is a distinct entity with its own timestamp. &
Binned: the timeline is divided into $T$ discrete slices. &
Exact: each De Bruijn node carries the exact timestamp of its contact(s). \\
\midrule
\textbf{Causal guarantee} &
Full: only forward-in-time edges exist. &
Partial: temporal encoding is in features; the GNN must learn causal rules. &
Full: only causally valid sequences exist as nodes; backward paths are absent. \\
\midrule
\textbf{Memory depth} &
1 hop (each event node has no memory of the path leading to it). &
All time steps (the full $T$-dim vector captures all past activations). &
$k-1$ hops explicitly encoded in the node identity. $k$ is a hyperparameter. \\
\midrule
\textbf{Key innovation} &
Lossless causal structure. First GNN approach to make time a topological property. &
Graph size unchanged. Temporal information as edge features; computationally efficient. &
Non-Markovian message passing. Sequences as nodes; the GNN reasons about paths, not hops. \\
\midrule
\textbf{Code location} &
\texttt{gnn/dag\_gnn.py},\newline \texttt{graph\_builder.py:\allowbreak build\_dag\_event\_graph} &
\texttt{gnn/backtracking\_network.py},\newline \texttt{graph\_builder.py:\allowbreak build\_temporal\_activation} &
\texttt{gnn/dbgnn.py},\newline \texttt{graph\_builder.py:\allowbreak build\_de\_bruijn\_graph} \\
\bottomrule
\end{tabular}
\caption{Comparison of the three temporal GNN architectures studied in this workbook. All three
solve the fundamental causal leakage problem of static GNNs, but through different mechanisms
with different trade-offs in graph size, memory depth, and computational cost.}
\label{tab:temporal_arch_comparison}
\end{table}

% =============================================================================
\section{Empirical Results from Qarkaxhija et al.}
% =============================================================================

The paper evaluates DBGNN on six time-series datasets: a synthetic ``temp-clusters'' benchmark
(designed to have structure exclusively in the causal topology) and five empirical contact
datasets (high-school proximity, hospital ward, student SMS messages, workplace contacts).

On the synthetic benchmark, all three time-aware methods (EVO, HONEM, DBGNN) achieve 100\%
balanced accuracy while static GCN achieves only 33.5\%. This confirms that the temporal
structure is necessary and sufficient for the task.

On empirical datasets, DBGNN achieves the best balanced accuracy on all five, with relative
improvements over the second-best method ranging from 1.55\% (workplace) to 28.16\% (hospital).
Table~\ref{tab:dbgnn_results} reproduces the key numbers.

\begin{table}[H]
\centering
\begin{tabular}{lrrrr}
\toprule
\textbf{Dataset} & \textbf{GCN} & \textbf{HONEM} & \textbf{DBGNN} & \textbf{Gain vs.\ 2nd best} \\
\midrule
temp-clusters   & 33.52 & 54.94 & \textbf{100.0} & 0.00\% \\
high-school-2011 & 50.06 & 54.24 & \textbf{64.4}  & +12.57\% \\
high-school-2012 & 58.03 & 53.16 & \textbf{65.8}  & +8.31\% \\
hospital        & 49.48 & 46.17 & \textbf{59.04} & +16.68\% \\
student-sms     & 57.33 & 50.43 & \textbf{60.6}  & +5.70\% \\
workplace       & 81.86 & 73.26 & \textbf{83.13} & +1.55\% \\
\bottomrule
\end{tabular}
\caption{Balanced accuracy (\%) for node classification on six dynamic graph datasets, comparing
static GCN, HONEM (a higher-order embedding baseline), and DBGNN. Results reproduced from
Qarkaxhija et al.\ (2022), Table~2. ``Gain'' is computed relative to the second-best method
on each dataset.}
\label{tab:dbgnn_results}
\end{table}

The hospital dataset shows the largest gain (+16.68\%). This is consistent with the domain:
hospital contact patterns are highly structured around professional roles (doctors, nurses,
patients, administrative staff), and the \emph{sequences} of contacts---who visits whom in
what order---are far more informative about role than any individual contact in isolation.

% =============================================================================
\section{Limitations and Open Questions}
% =============================================================================

\begin{keyinsight}{What DBGNN Does Not Solve}
Despite its theoretical appeal, the De Bruijn approach has several practical limitations that
are worth understanding before applying it:

\begin{enumerate}
    \item \textbf{Graph construction cost.} The De Bruijn graph must be built anew for every
          network and every time window. For a fixed network with static contact sequences, this
          is a one-time preprocessing cost. But for streaming data or online inference, it
          requires frequent recomputation.

    \item \textbf{Memory for high $k$.} The number of De Bruijn nodes grows rapidly with $k$
          (the number of valid $k$-step sequences can be exponential in $k$ for dense networks).
          In practice, the paper uses $k=2$ for all real datasets and finds that this is
          sufficient. Our implementation uses $k=1$ (first-order De Bruijn, equivalent to mapping
          individual directed contacts to nodes).

    \item \textbf{Source detection vs.\ node classification.} The paper was designed for a
          \emph{node classification} task (e.g., classify each person by their department).
          Adapting it to \emph{source detection} (identify the single epidemic seed node)
          requires the sentinel readout and susceptibility masking described in Section~6.
          These additions are specific to our epidemic application and are not in the original
          paper.

    \item \textbf{Order selection is unsupervised.} The statistical model selection for $k$
          is unsupervised: it uses only the topology of the data, not the labels. This is a
          strength (no label leakage) but also means the selected $k$ may not be optimal for
          the downstream classification task specifically. In principle, $k$ could be tuned
          jointly with other hyperparameters.
\end{enumerate}
\end{keyinsight}

% =============================================================================
\section{Chapter Summary}
% =============================================================================

This chapter introduced the De Bruijn Graph Neural Network (DBGNN), proposed by Qarkaxhija,
Perri, and Scholtes (2022).

The key ideas, in order:

\begin{enumerate}
    \item \textbf{The problem:} Standard GNNs on static graphs are Markovian. They allow
          messages to flow backward in time, creating causal leakage. Even time-aware approaches
          like the Backtracking Network remain partly Markovian because they aggregate
          individual contacts, not sequences.

    \item \textbf{The solution:} Replace the original graph with a \emph{De Bruijn graph} in
          which nodes are contact \emph{sequences} (causal walks) and edges connect sequences
          that share a valid overlap. On this graph, message passing is inherently
          non-Markovian: every node already ``remembers'' its full $k$-step history.

    \item \textbf{The sentinel trick:} Because De Bruijn nodes are events, not people, we add
          sentinel nodes $(v, v, t_{\text{end}})$ as readout anchors. After message passing,
          the sentinel's embedding summarises all causal evidence involving person $v$, and is
          used to score $v$ as a potential epidemic source.

    \item \textbf{Implementation:} The complete pipeline is implemented in
          \texttt{utils/make\_de\_bruijn\_graph.py} (graph construction),
          \texttt{gnn/graph\_builder.py} (tensor preparation), and
          \texttt{gnn/dbgnn.py} (the DBGNN model). The DBGNN model is approximately 178 lines
          of PyTorch, with the most critical logic in the \texttt{DBGNNLayer} scatter-add
          aggregation and the dual-projection feature initialisation.

    \item \textbf{Empirical evidence:} DBGNN outperforms static GNNs by large margins on
          temporally structured datasets, and outperforms time-aware baselines on all five
          empirical contact datasets evaluated in the paper.
\end{enumerate}

The next chapter examines the classical (non-learning) baselines for epidemic source detection,
providing a benchmark against which the GNN-based methods are compared.
